\chapter{Ejemplos}

\section{FFT con eje de frecuencia}

\begin{lstlisting}[language=Matlab]
	function [X,w] = contfft(x,T)
	% CONTFFT [X,w] = contfft(x,T)
	%  
	% Computes the Fourier transform of a continuous time signal
	% using the FFT. The input is the sampled continuous
	% time signal x and the sampling time T. The output is
	% the Fourier transform X(w) and the frequency vector w.
	%
	[n,m] = size(x);
	if n < m
	x = x';
	end
	Xn = fft(x);
	N = length(x);
	n = 0:N-1;
	n(1) = eps;
	X = (1-exp(-1j*2*pi*n/N))./(1j*2*pi*n/N/T).*Xn.';
	w = 2*pi*n/N/T;
\end{lstlisting}

\section{Filtro y FFT}

\begin{lstlisting}[language=Matlab]
	% Frecuencia de muestreo en muestras por segundo
	Fm  = 20000;  
	% La mitad de la frecuencia de muestreo
	Fmm = Fm/2;
	% Frecuencia en Hz a la que tiene que haber rp decibeles de atenuacion   
	wp  = 5000;
	% Atenuacion en dB a la frecuencia de corte   
	rp  = 3;
	% Frecuencia en Hz a la cual tiene que haber rs decibeles de atenuacion      
	ws  = 6000; 
	% Atenuacion minima en la banda de rechazo  
	rs  = 40;
	
	% buscamos el orden necesario
	[n1,wn1] = buttord(wp/Fmm, ws/Fmm, rp, rs);
	
	% calculamos el filtro
	[b1,a1] = butter(n1,wn1);
	
	% graficamos la respuesta en frecuencia del filtro
	figure(1); 
	clf; 
	[h1,w1] = freqz(b1,a1,512);
	figure(1);  
	clf;
	subplot(2,1,1); plot(w1*Fmm/pi,abs(h1));
	subplot(2,1,2); plot(w1*Fmm/pi,unwrap(angle(h1)));
	
	% graficamos la respuesta temporal
	t   = (1:30000)/Fm;
	s1  = sin(2*pi*t*1000);
	s2  = sin(2*pi*t*5000);
	s3  = sin(2*pi*t*6000);
	s   = s1 + 2*s2 + 1/4*s3;
	sf1 = filter(b1,a1,s);
	ns  = 20;
	
	S1  = fft(s,ns)/(ns/2);
	s2  = s1(20:1:(1024-1));
	S2  = fft(s2,ns);
	SF1 = fft(sf1,ns)/(ns/2);
	w1  = (0:(ns/2-1))/(ns/2)*Fmm;
	figure(2);
	clf;
	plot(w1,abs([S1(1:ns/2)' SF1(1:ns/2)']));
\end{lstlisting}


\section{Dise\~{n}o de filtros de tiempo discreto}

\subsection{Filtro Butterworth pasabajo}

\begin{lstlisting}[language=Matlab]
	% filtro Butterworth pasabajo de tiempo discreto
	
	Wm = 8500;  % frecuencia de muestreo
	Wmm = Wm/2; % mitad de la frecuencia de muestreo
	W1 = 1500;  % frecuencia superior banda de paso                     
	Wp = W1/Wmm; % frecuencia superior banda de paso normalizada 
	rp = 3;     % ripple banda de paso en dB
	W2 = 1.1*W1; % frecuencia inferior banda de rechazo 
	Ws = W2/Wmm; % frecuencia inferior banda de rechazo normalizada
	rs = 20;    % atenuacion banda de rechazo en dB 
	
	% usamos buttord para dimensionar el filtro
	[N,Wn] = buttord(Wp,Ws,rp,rs);
	
	% calculamos el filtro
	[b,a] = butter(N,Wn);
	
	% graficamos respuesta en frecuencia
	dbode(b,a,1/Wm);
\end{lstlisting}

\subsection{Filtro Butterworth pasabanda}

\begin{lstlisting}[language=Matlab]
	% FILTRO Butterworth pasabanda
	Wm  = 11000;    % frecuencia de muestreo
	Wmm = Wm/2;     % mitad de la frecuencia de muestreo
	W1  = 2800;     % frecuencia minima banda de paso
	W2  = 4000;     % frecuencia maxima banda de paso
	W0  = 0.9*W1;   % frecuencia maxima banda de rechazo inferior
	W3  = 1.1*W2;   % frecuencia minima banda de rechazo superior
	rp  = 3;        % ripple banda de paso
	rs  = 20;       % atenuacion banda de rechazo
	
	% usamos buttord para dimensionar el filtro
	[N,Wn] = buttord([W1/Wmm W2/Wmm], [W0/Wmm W3/Wmm], rp, rs);
	
	% calculamos el filtro
	[b,a] = butter(N,Wn);
	
	% graficamos respuesta en frecuencia
	dbode(b,a,1/Wm);
\end{lstlisting}

\subsection{Filtro el\'{\i}ptico pasabanda}

\begin{lstlisting}[language=Matlab]
	Wm = 100;               % frecuencia de muestreo
	t  = (1:100)/Wm;        % vector tiempo    
	s1 = sin(2*pi*t*5);     % componente de la senal a filtrar
	s2 = sin(2*pi*t*15);    % componente de la senal a filtrar
	s3 = sin(2*pi*t*30);    % componente de la senal a filtrar
	s  = s1 + s2 + s3;      % senal a filtrar
	plot(t,s);              % grafico de la senal a filtrar
	
	% Para disenar un filtro para mantener la sinusoide de 
	% 15 Hz y eliminar las sinusoides de 5 y 30 Hz, vamos a 
	% usar un filtro eliptico con una pasabanda de 10 a 20 Hz
	[b,a] = ellip(4,0.1,40,[10 20]*2/Wm); % diseno de filtro de tiempo discreto
	[H,w] = freqz(b,a,512);               % calculo de la respuesta en frecuencia
	plot(w*Wm/(2*pi),abs(H));             % grafico de la respuesta en frecuencia
	sf = filter(b,a,s);                   % filtrado
	plot(t,sf);                           % grafico de la senal filtrada
	xlabel('Tiempo (segundos)');          % rotulo x
	ylabel('Senal filtrada');             % rotulo y
	axis([0,1,-1,1]);                     % marco del grafico
	
	S = fft(s,512);                       % transformada rapida de Fourier
	SF = fft(sf,512);                     % transformada rapida de Fourier
	w = (0:255)/256*(Wm/2);               % vector de frecuencias para graficar el espectro
	
	% Observe que los picos de 5 y 30 Hz han sido eliminados
	plot(w, abs([S(1:256)' SF(1:256)']));
	xlabel('Frecuencia (Hz)');
	ylabel('Magnitud de la transformada de Fourier'); 
\end{lstlisting}

\subsection{Filtro el\'{\i}ptico suprimebanda}

\begin{lstlisting}[language=Matlab]
	Wm = 11000;
	Wmm = Wm/2;
	W1 = 2800;
	W2 = 4000;
	W0 = 1.1*W1;
	W3 = 0.9*W2;
	rp = 3;
	rs = 40;
	
	[N,Wn] = ellipord([W1/Wmm W2/Wmm], [W0/Wmm W3/Wmm], rp, rs);
	[b,a] = ellip(N, rp, rs, Wn, 'stop');
	dbode(b,a,1/Wm);
\end{lstlisting}

\begin{lstlisting}[language=Matlab]
	Wm = 8000;
	t  = (1:500)/Wm;
	s1 = sin(2*pi*t*500);
	s2 = sin(2*pi*t*1500);
	s3 = sin(2*pi*t*2500);
	s  = s1 + s2 + s3;
	
	figure;
	subplot(5,1,1); plot(t(400:500), s(400:500));
	subplot(5,1,2); plot(t(400:500), s1(400:500));
	subplot(5,1,3); plot(t(400:500), s2(400:500));
	subplot(5,1,4); plot(t(400:500), s3(400:500));
	
	dp = 0.1;
	rs = 40;
	n = 4;
	[b,a] = ellip(n, dp, rs, [1000,2000]*2/Wm);
	sf = filter(b,a,s);
	subplot(5,1,5); plot(t(400:500), sf(400:500));
\end{lstlisting}


\chapter{Ap\'{e}ndices y Tablas de Referencia}

\section{Conclusiones sobre filtros de tiempo discreto}

\begin{center}
	\renewcommand{\arraystretch}{1.3}
	\begin{tabular}{ll}
		\textbf{Ubicaci\'{o}n singularidad}                             & \textbf{Filtro}               \\ \hline
		Polos reales entre $z=0$ y $z=1$                                & pasabajo                      \\ 
		Ceros reales entre $z=-1$ y $z=0$                               & pasabajo                      \\ 
		Polos reales entre $z=-1$ y $z=0$                               & pasaalto                      \\ 
		Ceros reales entre $z=0$ y $z=1$                                & pasaalto                      \\ 
		Polos complejos conjugados dentro del c\'{\i}rculo unitario     & pasabanda                     \\ 
		Ceros complejos conjugados dentro del c\'{\i}rculo unitario     & suprimebanda 
	\end{tabular}
\end{center}

\begin{enumerate}
	\item Los filtros con polos complejos conjugados $\left| \alpha \right| e^{j\Omega }$ tienen frecuencias centrales que se incrementan cuando $\Omega $ se incrementa.
	\item Los filtros con ceros y polos complejos conjugados m\'{a}s cercanos al c\'{\i}rculo unidad tienen menor ancho de banda.
\end{enumerate}

\section{Transformaci\'{o}n de filtros de tiempo continuo a tiempo discreto}

Con las siguientes transformaciones de funciones de transferencias de tiempo continuo a tiempo discreto podemos convertir un filtro de tiempo continuo a tiempo discreto y viceversa.

\begin{center}
	Aproximaci\'{o}n de la integral \\ \vspace{0.2cm}
	\renewcommand{\arraystretch}{1.8}
	\begin{tabular}{ll}
		\textbf{Transformaci\'{o}n} & \textbf{F\'{o}rmula} \\ \hline
		Rectangular hacia adelante & $s \Leftrightarrow \frac{1-z^{-1}}{Tz^{-1}}$ \label{adelante} \\
		Rectangular hacia atr\'{a}s & $s \Leftrightarrow \frac{1-z^{-1}}{T}$ \label{atras}\\
		Trapezoidal (Bilineal) & $s \Leftrightarrow \frac{2}{T}\frac{1-z^{-1}}{1+z^{-1}}$ \label{trapezoidal}
	\end{tabular}
\end{center}

\vspace{0.5cm}

\begin{center}
	Aproximaci\'{o}n de la diferenciaci\'{o}n \\ \vspace{0.2cm}
	\renewcommand{\arraystretch}{1.8}
	\begin{tabular}{ll}
		\textbf{Transformaci\'{o}n} & \textbf{F\'{o}rmula} \\ \hline
		Diferencia hacia atr\'{a}s & $s \Leftrightarrow \frac{1-z^{-1}}{T}$
	\end{tabular}
\end{center}

\subsection{Transformaci\'{o}n bilineal}

Supongamos la ecuaci\'{o}n diferencial de primer orden 
\begin{equation}
	c_{1}\frac{dy\left( t\right) }{dt}+c_{0}y\left( t\right) =d_{0}x\left( t\right) ,
\end{equation}
con la funci\'{o}n de transferencia 
\begin{equation}
	H\left( s\right) =\frac{d_{0}}{c_{1}s+c_{0}}.
\end{equation}
Puede plantearse la siguiente relaci\'{o}n 
\begin{equation}
	y\left( t\right) -y\left( t_{0}\right) =\int_{t_{0}}^{t}\frac{dy\left( \tau \right) }{d\tau }d\tau ,
\end{equation}
y en particular 
\begin{equation}
	y\left( nT\right) -y\left( \left( n-1\right) T\right) =\int_{\left( n-1\right) T}^{nT}\frac{dy\left( \tau \right) }{d\tau }d\tau .
\end{equation}
Aproximando la integral por una regla trapezoidal, podemos escribir 
\begin{equation}
	y\left( nT\right) -y\left( \left( n-1\right) T\right) =\frac{T}{2}\left( 
	\left[ \frac{dy\left( \tau \right) }{d\tau }\right] _{\tau =nT}-\left[ 
	\frac{dy\left( \tau \right) }{d\tau }\right] _{\tau =\left( n-1\right) T}\right) ,
\end{equation}
y teniendo en cuenta que 
\begin{equation}
	c_{1}\left[ \frac{dy\left( \tau \right) }{d\tau }\right] _{\tau =nT}+c_{0}y\left( nT\right) =d_{0}x\left( nT\right) ,
\end{equation}
podemos finalmente obtener una ecuaci\'{o}n de diferencias asociada a la funci\'{o}n de transferencia 
\begin{equation}
	H\left( z\right) =\frac{d_{0}}{c_{1}\frac{2}{T}\frac{1-z^{-1}}{1+z^{-1}} + c_{0}}.
\end{equation}
Vemos entonces que podemos directamente usar la transformaci\'{o}n bilineal 
\begin{equation}
	s=\frac{2}{T}\frac{1-z^{-1}}{1+z^{-1}},
\end{equation}
y su inversa
\begin{equation}
	z=\frac{1+\frac{T}{2}s}{1-\frac{T}{2}s}.
\end{equation}

Aplicando la transformaci\'{o}n bilineal obtenemos el mapeo frecuencial:
\begin{equation}
	s_{p}=\left( -1\right) ^{\frac{1}{2N}}j\omega _{c}=\frac{2}{T}\frac{1-z_{p}^{-1}}{1+z_{p}^{-1}},
\end{equation}
o equivalentemente:
\begin{equation}
	z_{p}=\frac{1+\frac{T}{2} \left( -1\right)^{\frac{1}{2N}}j\omega _{c}}{1-\frac{T}{2}\left( -1\right) ^{\frac{1}{2N}}j\omega _{c}}.
\end{equation}

Por ejemplo, para un controlador PID de tiempo continuo tenemos una funci\'{o}n de transferencia:
\begin{equation}
	H(s) = K_{p} + \frac{K_{i}}{s} + K_{d} s
\end{equation}

Aplicando la transformaci\'{o}n rectangular hacia adelante (\ref{adelante}) obtenemos:
\begin{equation}
	H(z) = \frac{(K_{d} - K_{p} T + K_{i} T^{2}) + (K_{p} T - 2 K_{d}) z + K_{d} z^2}{T (z - 1)}.
\end{equation}

Aplicando la transformaci\'{o}n rectangular hacia atr\'{a}s (\ref{atras}) obtenemos:
\begin{equation}
	H(z) = \frac{K_{d} - (K_{p} T + 2 K_{d}) z + (K_{d} + K_{p} T + K_{i} T^{2}) z^2}{T (z - 1) z}.
\end{equation}

\section{Sumario de representaciones de se\~{n}ales en los dominios del tiempo y la frecuencia}

\begin{center}
	\renewcommand{\arraystretch}{2}
	\begin{tabular}{p{2cm} p{5.5cm} p{6cm}}
		\textbf{Categor\'{\i}a} & \textbf{Dominio del tiempo} & \textbf{Dominio de la frecuencia} \\ \hline
		\textbf{CTFT} & $x(t) =\int_{-\infty }^{\infty }\chi(f) e^{j2\pi ft}df$ & $\chi(f) =\int_{-\infty }^{\infty }x(t) e^{-j2\pi ft}dt$ \\ 
		\textbf{CTFS} & $x(t) =\sum_{m=-\infty }^{\infty }\chi[m] e^{\frac{j2\pi mt}{T}}$ & $\chi[m] =\frac{1}{T}\int_{-\frac{T}{2}}^{\frac{T}{2}}x(t) e^{-\frac{j2\pi mt}{T}}dt$ \\ 
		\textbf{DTFT} & $x[n] =\int_{-\frac{1}{2}}^{\frac{1}{2}}\bar{\chi}(\phi) e^{j2\pi \phi n}d\phi $ & $\bar{\chi}(\phi) =\sum_{m=-\infty }^{\infty }x[n] e^{-j2\pi \phi n}$ \\ 
		\textbf{DTFS} & $x[n] =\sum_{m=0}^{N-1}\bar{\chi}[m] e^{j2\pi \frac{mn}{N}}$ & $\bar{\chi}[m] =\frac{1}{N}\sum_{n=0}^{N-1}x[n] e^{-j2\pi \frac{mn}{N}}$
	\end{tabular}
\end{center}

\section{Propiedades comunes de los sistemas de los ejemplos}

Los sistemas definidos por las ecuaciones (\ref{sisUno}), (\ref{sisDos}), (\ref{sisTres}), y (\ref{sisCuatro}) tienen las siguientes propiedades en com\'{u}n:

\begin{enumerate}
	\item \textbf{Lineales}, porque las ecuaciones s\'{o}lo contienen operadores lineales.
	\item \textbf{Con memoria}, por las integrales, las derivadas y sumatorias.
	\item \textbf{Causales}, porque en las ecuaciones no aparecen valores futuros de la entrada, s\'{o}lo valores actuales o pasados.
\end{enumerate}

\section{Transformadas Laplace unilaterales}

\subsection{Funciones trigonom\'{e}tricas, primera parte}

\begin{center}
	\renewcommand{\arraystretch}{1.8}
	\begin{tabular}{ll}
		\textbf{$f(t)$} & \textbf{$F(s) = \mathcal{L}\{f(t)\}$} \\ \hline
		$\sin(\omega t)$ & $\frac{\omega }{s^{2}+\omega ^{2}}$ \\ 
		$\cos(\omega t)$ & $\frac{s}{s^{2}+\omega ^{2}}$ \\ 
		$\sin^{2}(\omega t)$ & $\frac{2\omega ^{2}}{s(s^{2}+4\omega ^{2})}$ \\ 
		$\cos^{2}(\omega t)$ & $\frac{s^{2}+2\omega ^{2}}{s(s^{2}+4\omega ^{2})}$ \\ 
		$t\sin(\omega t)$ & $\frac{2\omega s}{(s^{2}+\omega^{2})^{2}}$ \\ 
		$t\cos(\omega t)$ & $\frac{s^{2}-\omega ^{2}}{(s^{2}+\omega^{2})^{2}}$ \\ 
		$\frac{2(1-\cos(\omega t))}{t}$ & $\ln \left( \frac{s^{2}+\omega ^{2}}{s^{2}} \right)$ \\ 
		$\frac{\sin(\omega t)}{t}$ & $\arctan \left(\frac{\omega }{s}\right)$
	\end{tabular}
\end{center}

\subsection{Funciones trigonom\'{e}tricas, segunda parte} 

\begin{center}
	\renewcommand{\arraystretch}{1.8}
	\begin{tabular}{ll}
		\textbf{$f(t)$} & \textbf{$F(s) = \mathcal{L}\{f(t)\}$} \\ \hline
		$\sin(\omega t) + \omega t\cos(\omega t)$ & $\frac{2\omega s^{2}}{(s^{2}+\omega ^{2})^{2}}$ \\ 
		$\sin(\omega t) - \omega t\cos(\omega t)$ & $\frac{2\omega ^{3}}{(s^{2}+\omega ^{2})^{2}}$ \\ 
		$1-\cos(\omega t)$ & $\frac{\omega ^{2}}{s(s^{2}+\omega^{2})}$ \\ 
		$(\omega t) - \sin(\omega t)$ & $\frac{\omega ^{3}}{s^{2}(s^{2}+\omega ^{2})}$ \\ 
		$\frac{a\sin(bt) - b\sin(at)}{ab(a^{2}-b^{2})}$ & $\frac{1}{(s^{2}+a^{2})(s^{2}+b^{2})}$ \\ 
		$\frac{\cos(bt) - \cos(at)}{a^{2}-b^{2}}$ & $\frac{s}{(s^{2}+a^{2})(s^{2}+b^{2})}$ \\ 
		$\frac{\sin(at)\cos(bt)}{t}$ & $\frac{1}{2}\arctan \left(\frac{a+b}{s}\right) + \frac{1}{2}\arctan \left(\frac{a-b}{s}\right)$ \\ 
	\end{tabular}
\end{center}

\subsection{Funciones hiperb\'{o}licas, primera parte}

\begin{center}
	\renewcommand{\arraystretch}{1.8}
	\begin{tabular}{ll}
		\textbf{$f(t)$} & \textbf{$F(s) = \mathcal{L}\{f(t)\}$} \\ \hline
		$\sinh(kt)$ & $\frac{k}{s^{2}-k^{2}}$ \\ 
		$\cosh(kt)$ & $\frac{s}{s^{2}-k^{2}}$ \\ 
		$\sinh^{2}(kt)$ & $\frac{2k^{2}}{s(s^{2}-4k^{2})}$ \\ 
		$\cosh^{2}(kt)$ & $\frac{s^{2}-2k^{2}}{s(s^{2}-4k^{2})}$
	\end{tabular}
\end{center}

\subsection{Funciones hiperb\'{o}licas, segunda parte}

\begin{center}
	\renewcommand{\arraystretch}{1.8}
	\begin{tabular}{ll}
		\textbf{$f(t)$} & \textbf{$F(s) = \mathcal{L}\{f(t)\}$} \\ \hline
		$t\sinh(kt)$ & $\frac{2ks}{(s^{2}-k^{2})^{2}}$ \\ 
		$t\cosh(kt)$ & $\frac{s^{2}+k^{2}}{(s^{2}-k^{2})^{2}}$ \\ 
		$\frac{2(1-\cosh(kt))}{t}$ & $\ln \left( \frac{s^{2}-k^{2}}{s^{2}} \right)$ \\ 
	\end{tabular}
\end{center}

\subsection{Funciones exponenciales}

\begin{center}
	\renewcommand{\arraystretch}{1.8}
	\begin{tabular}{ll}
		\textbf{$f(t)$} & \textbf{$F(s) = \mathcal{L}\{f(t)\}$} \\ \hline
		$e^{at}$ & $\frac{1}{s-a}$ \\ 
		$t e^{at}$ & $\frac{1}{(s-a)^{2}}$ \\ 
		$t^n e^{at}$ & $\frac{n!}{(s-a)^{n+1}}, \quad n\geq 0$ \\ 
		$\frac{e^{bt}-e^{at}}{t}$ & $\ln \left( \frac{s-a}{s-b} \right)$ \\
		$\frac{1}{\sqrt{\pi t}}e^{-a^{2}/4t}$ & $\frac{e^{-a\sqrt{s}}}{\sqrt{s}}$ \\ 
		$\frac{a}{2\sqrt{\pi t^{3}}}e^{-a^{2}/4t}$ & $e^{-a\sqrt{s}}$ \\ 
		$\frac{e^{at}-e^{bt}}{a-b}$ & $\frac{1}{(s-a)(s-b)}$ \\ 
		$\frac{ae^{at}-be^{bt}}{a-b}$ & $\frac{s}{(s-a)(s-b)}$
	\end{tabular}
\end{center}


\section{Libros}
\begin{enumerate}
	
	\item Boaz Porat, 
	\textit{A Course in Digital Signal Processing}, 
	Wiley, 1997. 
	
	\item James McClellan et al., 
	\textit{Computer-Based Exercises for Signal Processing Using MATLAB 5},
	Prentice Hall, 1998.
	
	\item Ronald N. Bracewell, 
	\textit{The Fourier Transform and Its Applications},
	McGraw-Hill: New York, 1986.
	
	\item Roland E. Thomas and Albert J. Rosa, 
	\textit{The Analysis and Design of Linear Circuits},
	Prentice Hall, New Jersey, 1994.
	
\end{enumerate}